\usepackage[utf8]{inputenc}
\usepackage[T1]{fontenc}
% The report is in English (British, Oxford -ize spelling); babel takes the
% main language from the last option.  Swedish stays loaded for the few
% Swedish words the report quotes (course terms, feedback headings).
\usepackage[swedish,english,british]{babel}
\usepackage{booktabs}
\usepackage{array}
\usepackage{longtable}
\usepackage{tabularx}

\usepackage[
  natbib,
  style=verbose,
  citestyle=verbose,
  singletitle=false,
  maxbibnames=99,
]{biblatex}
\addbibresource{ltnotes.bib}

\usepackage[all]{foreign}
\renewcommand{\foreignfullfont}{}
\renewcommand{\foreignabbrfont}{}

\usepackage{import}
% beamer loads graphicx itself; memoir does not, and a deck with figures
% needs it in both jobs.
\usepackage{graphicx}

% The bibliography has long URLs and DOIs; let biblatex break them.
\setcounter{biburlnumpenalty}{7000}
\setcounter{biburlucpenalty}{7000}
\setcounter{biburllcpenalty}{7000}

\usepackage[strict]{csquotes}
\usepackage[single]{acro}

\usepackage{subcaption}

% The workflow figure is drawn with TikZ, so its text is set with the
% document at the size around it.
\usepackage{tikz}
\usetikzlibrary{positioning,arrows.meta}

\usepackage{xparse}

\let\email\texttt

% minted v3+ (TeX Live 2024+) auto-detects the output directory; the old
% [outputdir=...] package option is an error there.
\usepackage{minted}
\setminted{autogobble,fontsize=\footnotesize}

% PythonTeX runs the example programs and embeds their real terminal
% sessions (didactic's \runpython[transcript]), so the deck never
% hand-copies its own output: pythontex runs from ltxobj (the outputdir),
% and workingdir=.. moves it up to the deck directory so examples/...
% resolves.
\usepackage[makestderr]{pythontex}
\setpythontexoutputdir{.}
\setpythontexworkingdir{..}

\usepackage{amsmath}
\usepackage{amssymb}
\usepackage{mathtools}
\usepackage{amsthm}
\usepackage{thmtools,thm-restate}
\usepackage[unq]{unique}

\usepackage[binary-units,parse-numbers=false]{siunitx}

\makeatletter
\@ifclassloaded{beamer}
  {\relax}
  {\usepackage[colorlinks]{hyperref}}
\makeatother
% cleveref reads its language from its own package options, not babel's.
% English decks capitalise mid-sentence references (Figure 3, Appendix B).
\usepackage[english,capitalise]{cleveref}
\AtBeginDocument{%
  % Figures are "Figure 3" in the running text, not "Fig. 3".
  \crefname{figure}{Figure}{Figures}%
  \Crefname{figure}{Figure}{Figures}%
  % The prose refers back to numbered lines of a program listing.
  \crefname{enumi}{line}{lines}%
  \Crefname{enumi}{Line}{Lines}%
}

\usepackage[marginparmargin=outer]{didactic}
% didactic sets footnotes --- and, under the verbose citation style, every
% \autocite --- as margin notes, and LaTeX allocates \marginpar boxes from
% the same pool of eighteen boxes as floats.  A chapter with this many
% margin references empties the pool and dies with "Too many unprocessed
% floats" pages later, in a place that has nothing to do with the cause.
\extrafloats{200}
% memoir places a side caption with \rlap, so it takes part in no
% collision avoidance at all and silently overprints an \ltnote that
% happens to sit at the same height.  Aligning the caption with the BOTTOM
% of the figure moves it below the notes, which anchor where they stand in
% the body text.  (\setsidecappos is memoir's; beamer does not define it.)
\ifdefined\setsidecappos\setsidecappos{b}\fi
\ProvideSemanticEnv{lo}{Learning Objective}
  [style=definition,numbered=yes]
  {LO}{LO}
  {Learning objective}{Learning objectives}
\ProvideTranslation{british}{Learning Objective}{Learning Objective}
\ProvideTranslation{british}{LO}{LO}
\ProvideTranslation{british}{Learning objective}{Learning objective}
\ProvideTranslation{british}{Learning objectives}{Learning objectives}
\ProvideTranslation{swedish}{Learning Objective}{Lärandemål}
\ProvideTranslation{swedish}{LO}{lm}
\ProvideTranslation{swedish}{Learning objective}{Lärandemål}
\ProvideTranslation{swedish}{Learning objectives}{Lärandemål}

% Paper titles in the generated hit-list tables can carry Greek letters;
% scholar's table export sets them in math mode itself, and this mapping
% stays as a guard for the same characters typed anywhere else.
\DeclareUnicodeCharacter{03C0}{\ensuremath{\pi}}
\DeclareUnicodeCharacter{03BB}{\ensuremath{\lambda}}
\DeclareUnicodeCharacter{03B1}{\ensuremath{\alpha}}
\DeclareUnicodeCharacter{03B2}{\ensuremath{\beta}}
\DeclareUnicodeCharacter{03BC}{\ensuremath{\mu}}
% A Vietnamese title in a hit list has u with a horn, which T1 lacks; it
% is set as a plain u.
\DeclareUnicodeCharacter{01B0}{u}

% ------------------------------------------------------------------ figures
% Pictures and tables that appear both on a slide and in a float of the
% notes are defined once here, so the two cannot drift apart.  (A
% \newcommand in contents.nw would sit outside any frame, where beamer's
% \mode* skips it in the slides job.)

% The workflow of the rewrite, ten steps in two rows read as a snake.
\tikzset{
  wfstep/.style={
    draw, rounded corners, align=center, text width=1.7cm,
    minimum height=0.95cm, inner sep=2pt, font=\footnotesize,
  },
  wfarrow/.style={-{Stealth[length=4pt]}, thick},
}
\newcommand{\workflowfigure}{%
  \begin{tikzpicture}[node distance=3mm and 3mm]
    \node[wfstep] (s1) {1.\ Export and verify};
    \node[wfstep, right=of s1] (s2) {2.\ Decide the draw};
    \node[wfstep, right=of s2] (s3) {3.\ Style guide};
    \node[wfstep, right=of s3] (s4) {4.\ Author decides};
    \node[wfstep, right=of s4] (s5) {5.\ Seven rewriters};
    \node[wfstep, below=of s5] (s6) {6.\ Three reviewers};
    \node[wfstep, left=of s6] (s7) {7.\ Fixes};
    \node[wfstep, left=of s7] (s8) {8.\ Feedback};
    \node[wfstep, left=of s8] (s9) {9.\ Draw program};
    \node[wfstep, left=of s9] (s10) {10.\ Example exam};
    \foreach \a/\b in {s1/s2, s2/s3, s3/s4, s4/s5, s5/s6,
                       s6/s7, s7/s8, s8/s9, s9/s10}
      \draw[wfarrow] (\a) -- (\b);
  \end{tikzpicture}%
}

% What each step guards against, as table rows, in two halves so that a
% slide can show one half and the notes' table both.
\newcommand{\workflowtablehead}{%
  \toprule
  Step & What it guards against \\
  \midrule
}
\newcommand{\workflowrowsA}{%
  1.\ Export, convert, round trip through a sandbox
    & a conversion that silently changes a question \\
  2.\ Decide how to draw
    & a draw the platform cannot do; a retake meeting the same question \\
  3.\ Style guide from the course and its live coding
    & items in a style the students never met, or whose style gives the
      key away \\
  4.\ The author decides the open points
    & agents settling what is the teacher's to settle \\
  5.\ Seven rewriters, every program run
    & wrong keys, crashes, several valid orders, cues \\
}
\newcommand{\workflowrowsB}{%
  6.\ Three independent reviewers
    & what the rewriters cannot see in their own work \\
  7.\ Fixes by the original rewriters
    & a fix that breaks something else in the bank \\
  8.\ Feedback on every question and option
    & feedback that is wrong, names the key, or points at material that
      does not exist \\
  9.\ A literate, reviewed drawing program
    & a draw nobody can reproduce; exam content in the public repository;
      a retake meeting a question again \\
  10.\ An example exam in the sandbox
    & a difference between the draw and what Canvas shows \\
}
% Column types for that table: a ragged step column and a ragged X column.
\newcolumntype{W}{>{\raggedright\arraybackslash}p{0.36\linewidth}}
\newcolumntype{R}{>{\raggedright\arraybackslash}X}

% The number of questions per bank, before and after the rewrite.  Bank
% L.n is the n-th bank of learning objective L.
\newcommand{\banktable}{%
  \begin{tabular}{@{}lrrrr@{}}
    \toprule
    Objective & Bank 1 & Bank 2 & Bank 3 & All three \\
    \midrule
    1 & $18\to12$ & $9\to9$   & $8\to8$   & $35\to29$ \\
    2 & $14\to12$ & $11\to10$ & $12\to9$  & $37\to31$ \\
    3 & $11\to10$ & $10\to9$  & $8\to8$   & $29\to27$ \\
    4 & $13\to10$ & $10\to10$ & $10\to10$ & $33\to30$ \\
    5 & $7\to9$   & $10\to9$  & $11\to9$  & $28\to27$ \\
    6 & $8\to10$  & $12\to10$ & $12\to9$  & $32\to29$ \\
    7 & $9\to9$   & $7\to8$   & $8\to8$   & $24\to25$ \\
    \midrule
    All & & & & $218\to198$ \\
    \bottomrule
  \end{tabular}%
}
